% META: {"id":"1SPE-DERLOCAL-EX-035","chapitre":"1SPE-DERIVATION-LOCAL","type_objet":"exercice","capacites":["1SPE-DERIVATION-LOCAL-C3"],"capacites_codes":["C3"],"methodes":["M1"],"parcours":1,"competences":["representer"],"duree_min":8,"mode_creation":"ex_nihilo","sources_inspiration":[],"fichier_tex":"chapitres/1SPE-DERIVATION-LOCAL/exercices/1SPE-DERLOCAL-EX-035.tex","corrige_tex":"chapitres/1SPE-DERIVATION-LOCAL/corriges/1SPE-DERLOCAL-CO-035.tex","status":"approved","origin":"generated","status_history":["generated","audited","accepted","approved"],"release_acceptance":"RELEASE_OWNER_BATCH_ACCEPTANCE","acceptance_closure_digest":"sha256:9b3ccf9a81c5520fbb7e03b7b2d3e7bf2057a02834b6d908bf3be5f49f3b553f","acceptance_receipt_digest":"sha256:4fad86f0282e0fa4e0419cca1ad6379ae7af5a280c04a4a90880f90a1c044242"}
% BEGIN-VERIFY
% # tangente en x=2 : y = f'(2)(x-2) + f(2) = -1*(x-2) + 3 = -x + 5
% # point (2,3): y = -2+5 = 3 OK
% # point (3,2): y = -3+5 = 2 OK
% assert -1*(2-2) + 3 == 3
% assert -1*(3-2) + 3 == 2
% END-VERIFY
\begin{exercice}{1SPE-DERLOCAL-EX-035}{1}{8}
On sait que $f(2)=3$ et $f'(2)=-1$. Tracer, sur un repère, la tangente à la courbe de $f$ au point d'abscisse $2$.

On précisera les coordonnées d'au moins deux points de cette tangente.
\end{exercice}
